\documentclass[
    language=english,
    doctype=manual,
    institution=none,
    titlestyle=book,
    oneside,
]{../../omnilatex}

\RequirePackage{config/document-settings}

\addbibresource{../../bib/bibliography.bib}

\begin{document}

\maketitle

\tableofcontents

%=============================================================================
\chapter{Introduction}
\label{sec:introduction}
%=============================================================================

DataPipe is a streaming ETL platform for real-time data integration. This
user guide walks you through installation, configuration, everyday usage,
and advanced features. Whether you are evaluating DataPipe for a new project
or integrating it into an existing data stack, this guide provides the
information you need to get productive quickly.

\section{Who Should Read This Guide}

This guide is written for data engineers, platform operators, and
technical leads who need to move data between systems in real time.
Basic familiarity with the command line and REST APIs is assumed. No
prior experience with Apache Kafka or Flink is required---the guide
covers the necessary background in~\cref{sec:concepts}.

\section{Key Concepts}
\label{sec:concepts}

\begin{description}
    \item[Pipeline] A directed acyclic graph (DAG) of source, transform,
        and sink stages that describes the flow of data from origin to
        destination.
    \item[Connector] A pluggable module that reads from or writes to an
        external system (e.g.\ PostgreSQL, S3, Kafka topic).
    \item[Transform] An inline function applied to each record as it
        passes through a pipeline stage. Transforms are composable and
        stateless by default.
    \item[Window] A time-based or count-based grouping of records used
        for aggregations and sessionization.
    \item[Backpressure] The mechanism by which DataPipe slows ingestion
        when downstream consumers cannot keep up.
\end{description}

\section{System Requirements}

\begin{table}[htbp]
    \centering
    \caption{Minimum and recommended system requirements.}
    \label{tab:requirements}
    \begin{tabular}{@{}lll@{}}
        \toprule
        Component & Minimum & Recommended \\
        \midrule
        CPU cores & 4 & 8 \\
        RAM & 8\,GB & 32\,GB \\
        Disk (SSD) & 50\,GB & 200\,GB \\
        OS & Ubuntu 20.04+ & Ubuntu 22.04 LTS \\
        Java & 17 & 21 LTS \\
        \bottomrule
    \end{tabular}
\end{table}

%=============================================================================
\chapter{Installation}
\label{sec:installation}
%=============================================================================

\section{Prerequisites}

Before installing DataPipe, ensure that the following software is present
on your system:

\begin{enumerate}
    \item A supported operating system (see \cref{tab:requirements}).
    \item Java Development Kit (JDK) 17 or later. Verify with:
    \begin{lstlisting}[language=bash]
java -version
    \end{lstlisting}
    \item Docker 20.10+ and Docker Compose v2 (for the containerised
        deployment path).
    \item curl or wget for downloading release archives.
\end{enumerate}

\section{Quick Install via Script}

The fastest way to install DataPipe on a single node is the bootstrap
script. It downloads the latest release, verifies the checksum, and
installs the \texttt{dp} CLI into \texttt{/usr/local/bin}:

\begin{lstlisting}[language=bash, caption={Bootstrap installation script.}]
curl -fsSL https://get.datapipe.dev | bash
dp --version
\end{lstlisting}

\section{Docker Compose Deployment}

For development and evaluation, the included Docker Compose file starts a
three-node cluster with Kafka, Zookeeper, and the DataPipe runtime:

\begin{lstlisting}[language=bash, caption={Start the development cluster.}]
cd /opt/datapipe
docker compose up -d
docker compose logs -f datapipe
\end{lstlisting}

\section{Manual Installation}

For production deployments where Docker is not available:

\begin{lstlisting}[language=bash, caption={Manual tarball installation.}]
wget https://releases.datapipe.dev/3.2.1/datapipe-3.2.1.tar.gz
tar xzf datapipe-3.2.1.tar.gz -C /opt/
ln -s /opt/datapipe-3.2.1 /opt/datapipe
export PATH="/opt/datapipe/bin:$PATH"
\end{lstlisting}

Add the \texttt{PATH} export to your \texttt{\textasciitilde/.bashrc} or
\texttt{/etc/profile.d/datapipe.sh} for persistence.

%=============================================================================
\chapter{Configuration}
\label{sec:configuration}
%=============================================================================

\section{Configuration File Format}

DataPipe uses YAML configuration files. The primary configuration file is
\texttt{datapipe.yaml}, located in the installation directory or specified
via the \texttt{DP\_CONFIG} environment variable.

\begin{lstlisting}[language=yaml, caption={Example datapipe.yaml configuration.}]
version: "3.2"
cluster:
  name: production-cluster
  id: dc-prod-01
  seeds:
    - 10.0.1.10:9092
    - 10.0.1.11:9092
  replication_factor: 3
storage:
  type: local
  path: /data/datapipe
  retention_ms: 604800000  # 7 days
logging:
  level: INFO
  file: /var/log/datapipe/datapipe.log
  max_size_mb: 500
  compress: true
\end{lstlisting}

\section{Environment Variables}

Key environment variables override configuration file values:

\begin{table}[htbp]
    \centering
    \caption{DataPipe environment variables.}
    \label{tab:env-vars}
    \begin{tabular}{@{}lp{9cm}@{}}
        \toprule
        Variable & Description \\
        \midrule
        \texttt{DP\_CONFIG} & Path to the primary YAML configuration file. \\
        \texttt{DP\_HOME} & Installation root directory. \\
        \texttt{DP\_KAFKA\_BROKERS} & Comma-separated list of Kafka broker addresses. \\
        \texttt{DP\_LOG\_LEVEL} & Logging level (\texttt{DEBUG}, \texttt{INFO}, \texttt{WARN}, \texttt{ERROR}). \\
        \texttt{DP\_JVM\_HEAP} & JVM heap size (e.g.\ \texttt{8g}). \\
        \bottomrule
    \end{tabular}
\end{table}

\section{Pipeline Configuration}

A pipeline is defined in a separate YAML file and registered with the
DataPipe runtime. Each pipeline declares one or more stages:

\begin{lstlisting}[language=yaml, caption={Example pipeline definition.}]
pipeline:
  name: user-events-to-warehouse
  description: Ingest user events and load into the data warehouse
  stages:
    - name: source
      type: kafka-source
      config:
        topic: user-events
        consumer_group: etl-warehouse
        bootstrap_servers: ${DP_KAFKA_BROKERS}

    - name: transform
      type: filter
      config:
        predicate: "event_type IN ('click', 'purchase')"

    - name: sink
      type: postgres-sink
      config:
        table: analytics.user_events
        batch_size: 5000
        flush_interval_ms: 10000
\end{lstlisting}

%=============================================================================
\chapter{Using DataPipe}
\label{sec:usage}
%=============================================================================

\section{The \texttt{dp} Command-Line Interface}

The \texttt{dp} CLI is the primary interface for managing DataPipe.
Common commands include:

\begin{lstlisting}[language=bash, caption={Common dp CLI commands.}]
# List all registered pipelines
dp pipeline list

# Start a pipeline
dp pipeline start user-events-to-warehouse

# Check pipeline status
dp pipeline status user-events-to-warehouse

# View pipeline metrics
dp metrics user-events-to-warehouse --last 1h

# Stop a pipeline gracefully
dp pipeline stop user-events-to-warehouse
\end{lstlisting}

\section{Monitoring with the Web Dashboard}

DataPipe ships with a built-in web dashboard on port 8080. Navigate to
\texttt{http://localhost:8080} after starting the runtime to access:

\begin{itemize}
    \item Real-time throughput graphs per pipeline.
    \item Consumer lag monitoring.
    \item Error rate and retry statistics.
    \item Pipeline topology visualisation.
\end{itemize}

\section{REST API}

The same operations available through the CLI are also exposed via a REST
API:

\begin{lstlisting}[language=bash, caption={Example REST API calls.}]
# List pipelines
curl -s http://localhost:8080/api/v1/pipelines | jq .

# Get pipeline status
curl -s http://localhost:8080/api/v1/pipelines/user-events-to-warehouse | jq .

# Trigger a manual restart
curl -X POST http://localhost:8080/api/v1/pipelines/user-events-to-warehouse/restart
\end{lstlisting}

%=============================================================================
\chapter{Advanced Features}
\label{sec:advanced}
%=============================================================================

\section{Exactly-Once Semantics}

DataPipe supports exactly-once delivery for Kafka sources and supported
sinks. Enable it in the pipeline configuration:

\begin{lstlisting}[language=yaml, caption={Exactly-once pipeline configuration.}]
pipeline:
  name: exactly-once-example
  delivery_guarantee: exactly-once
  stages:
    - name: source
      type: kafka-source
      config:
        topic: financial-transactions
        isolation_level: read_committed
    - name: sink
      type: postgres-sink
      config:
        table: ledger.transactions
        idempotent: true
\end{lstlisting}

\section{Custom Transforms}

Write custom transforms in Java or Python. The transform must implement
the \texttt{Transform} interface:

\begin{lstlisting}[language=java, caption={Custom Java transform.}]
public class CurrencyConverter implements Transform {
    private final String sourceCurrency;
    private final String targetCurrency;
    private final ExchangeRateService rates;

    public CurrencyConverter(Config config) {
        this.sourceCurrency = config.getString("source_currency");
        this.targetCurrency = config.getString("target_currency");
        this.rates = new ExchangeRateService(config);
    }

    @Override
    public Record apply(Record record) {
        double amount = record.getDouble("amount");
        double rate = rates.getRate(sourceCurrency, targetCurrency);
        record.set("amount", amount * rate);
        record.set("currency", targetCurrency);
        return record;
    }
}
\end{lstlisting}

\section{Windowed Aggregations}

DataPipe supports tumbling, sliding, session, and global windows:

\begin{lstlisting}[language=yaml, caption={Windowed aggregation pipeline.}]
pipeline:
  name: hourly-revenue
  stages:
    - name: source
      type: kafka-source
      config:
        topic: purchases
    - name: window
      type: tumbling-window
      config:
        size_ms: 3600000  # 1 hour
        trigger: on_window_close
    - name: aggregate
      type: aggregate
      config:
        group_by: ["region", "product_category"]
        metrics:
          - field: amount
            function: sum
            output_field: total_revenue
          - field: order_id
            function: count
            output_field: order_count
    - name: sink
      type: postgres-sink
      config:
        table: analytics.hourly_revenue
\end{lstlisting}

%=============================================================================
\chapter{Troubleshooting}
\label{sec:troubleshooting}
%=============================================================================

\section{Common Issues}

\begin{description}
    \item[Pipeline fails to start]
        Check that all required configuration values are set. Run
        \texttt{dp pipeline validate <name>} to identify missing fields.
        Ensure the Kafka brokers are reachable from the DataPipe node.

    \item[High consumer lag]
        Consumer lag indicates the pipeline cannot keep up with the source
        throughput. Consider increasing the number of consumer instances,
        optimising transforms, or scaling the sink batch size.

    \item[Out-of-memory errors]
        Increase the JVM heap via \texttt{DP\_JVM\_HEAP} or the
        \texttt{storage.jvm\_heap} configuration key. If the issue
        persists, profile the pipeline to identify memory-intensive
        transforms.

    \item[Connection refused to Kafka]
        Verify broker addresses in the configuration, check firewall
        rules, and ensure Kafka is running. Use
        \texttt{dp cluster health} to diagnose connectivity.
\end{description}

\section{Log Analysis}

DataPipe logs to \texttt{/var/log/datapipe/datapipe.log} by default. To
enable debug logging temporarily:

\begin{lstlisting}[language=bash]
dp config set logging.level DEBUG
dp pipeline restart <name>
\end{lstlisting}

\section{Diagnostic Commands}

\begin{lstlisting}[language=bash, caption={Diagnostic and health-check commands.}]
dp cluster health
dp pipeline validate user-events-to-warehouse
dp pipeline status user-events-to-warehouse --verbose
dp metrics user-events-to-warehouse --last 24h --format csv > metrics.csv
dp logs --tail 100 --level error
\end{lstlisting}

%=============================================================================
\chapter{Frequently Asked Questions}
\label{sec:faq}
%=============================================================================

\begin{description}
    \item[Can DataPipe run on Kubernetes?]
        Yes. Helm charts are available at
        \texttt{https://charts.datapipe.dev}. See the Kubernetes
        deployment guide for cluster sizing recommendations.

    \item[What databases are supported as sinks?]
        PostgreSQL, MySQL, MariaDB, ClickHouse, BigQuery, Snowflake, and
        Redshift are officially supported. Community connectors exist for
        additional databases.

    \item[How do I upgrade DataPipe?]
        Run the bootstrap script again---it will detect the existing
        installation and perform an in-place upgrade. Always back up your
        configuration before upgrading.

    \item[Is there a managed cloud offering?]
        DataPipe Cloud is available at \texttt{https://cloud.datapipe.dev}
        with a free tier for evaluation workloads.

    \item[Can I use DataPipe without Kafka?]
        Yes. DataPipe includes a file-based connector for development and
        testing environments. Production deployments are recommended to use
        Kafka for durability and scalability.
\end{description}

\printbibliography[heading=bibintoc, title={References}]

\end{document}
